\documentclass[singleLetter,times]{sagej}

\usepackage{generic}
\usepackage{graphicx}
\usepackage{amsmath,amssymb,amsfonts}
\usepackage{cite}
\usepackage{textcomp}

% Match the compact page margins shown in the reference PDF.
\geometry{left=0.75in,right=0.75in,top=0.75in,bottom=0.75in}

\begin{document}

\title{Short-Trajectory Artifact-Reversal Diffusion for CCTA Calcium Deblooming}

\author{Mahmud Wasif Nafee\affil{1}, Muhan Shao\affil{2}, Christopher Wiedeman\affil{2}, Lin Fu\affil{2}, Bruno De Man\affil{2}, and Ge Wang\affil{1}}
\affiliation{\affilnum{1}Biomedical Imaging Center, Center for Computational Innovations, Center for Biotechnology and Interdisciplinary Studies, Department of Biomedical Engineering, School of Engineering, Rensselaer Polytechnic Institute, Troy, NY, USA\\
\affilnum{2}GE HealthCare, Technology and Innovation Center, Niskayuna, NY, USA}

\corrauth{Ge Wang, Biomedical Imaging Center, Center for Computational Innovations, Center for Biotechnology and Interdisciplinary Studies, Department of Biomedical Engineering, School of Engineering, Rensselaer Polytechnic Institute, Troy, NY, USA}
\email{wangg6@rpi.edu}

\begin{abstract}
\textbf{BACKGROUND:} Calcium blooming in coronary CT angiography (CCTA) enlarges calcified plaques, obscures the coronary lumen, and complicates stenosis assessment. Short-trajectory image restoration methods offer an efficient alternative to long-trajectory diffusion generation by progressively transforming a corrupted image toward a task-specific target.
\textbf{OBJECTIVE:} To develop a short-trajectory image restoration method for CCTA calcium deblooming that suppresses blooming artifacts while preserving the calcium core and recovering adjacent lumen structure.
\textbf{METHODS:} We propose Short-Trajectory Artifact-Reversal (STAR) Diffusion, which maps a blooming-corrupted patch toward a debloomed target through a localized short trajectory. We additionally introduce Soft Annular Boundary Attention (SABA) to reduce attention leakage from the blooming-prone annulus into the calcium core and support boundary-consistent restoration.
\textbf{RESULTS:} In 232 held-out synthetic test volumes, STAR Diffusion achieved a Dice score of 0.9250, Core RMSE of 134.50 HU, and Annulus RMSE of 90.24 HU, with significant gains in Dice and Annulus RMSE over classical, GAN-based, and latent diffusion baselines ($p<0.05$). Ablations showed complementary benefits from spatial prior conditioning and SABA, and the proposed method ran substantially faster than long-trajectory latent diffusion baselines.
\textbf{CONCLUSIONS:} Short-trajectory, spatially guided diffusion with boundary-aware attention can debloom CCTA while preserving calcium morphology and adjacent lumen structure, supporting improved coronary stenosis assessment in time- and resource-constrained settings.
\end{abstract}

\keywords{Diffusion models, coronary CT angiography, calcium blooming, image restoration, reconstruction artifacts}

\maketitle

\section{Introduction}
\label{sec:introduction}

Calcium blooming is an imaging artifact in \textbf{coronary CT angiography (CCTA)} where dense calcium plaques appear larger than their true size; examples of blooming-corrupted patches are shown in the first column of Fig.~\ref{fig:bloom_debloom_grid}. Because CCTA is used to assess clinically significant coronary stenosis, this plaque-size overestimation can obscure the coronary lumen and lead to false-positive stenosis assessment. Blooming has been linked to partial volume averaging, cardiac motion, and beam hardening~\cite{pack_blooming}, which blur the calcium-lumen boundary and complicate CCTA interpretation. As a result, calcium deblooming, or blooming artifact reduction, is needed to suppress artifact-induced plaque enlargement while preserving the true calcified plaque signal and improving lumen visualization.
Compared with other CT enhancement tasks, such as denoising and super-resolution, learning-based calcium deblooming is more challenging because the artifact is spatially coupled with the underlying plaque rather than being a separate or globally distributed degradation. As shown in Fig.~\ref{fig:bloom_debloom_grid}, blooming can simultaneously obscure part of the coronary lumen and exaggerate the apparent size of calcified plaques. Therefore, an effective deblooming model must address two coupled objectives: \textbf{lumen recovery}, where hidden or obscured vessel regions are reconstructed in a structurally consistent manner, and \textbf{calcium preservation}, where the true calcium core is separated from the surrounding blooming artifact so that the calcified plaque signal is preserved while the blooming halo is reduced.

\begin{figure}[t]
    \centering
    \setlength{\tabcolsep}{3pt}
    \renewcommand{\arraystretch}{0.9}

    \begin{tabular}{@{}lcccc@{}}
        \shortstack[l]{\small\textbf{with}\\\small\textbf{blooming}\\\small\textbf{artifact}} &
        \includegraphics[width=0.17\linewidth]{../tmbe_deblooming-short-traj-latex-main/images_1/03_ordinary_CaP_40_minus_artery_100mA_10_42_bloom.png} &
        \includegraphics[width=0.17\linewidth]{../tmbe_deblooming-short-traj-latex-main/images_1/32_ordinary_CaP_50_minus_artery_300mA_12_37_bloom.png} &
        \includegraphics[width=0.17\linewidth]{../tmbe_deblooming-short-traj-latex-main/images_1/35_ordinary_CaP_50_minus_artery_100mA_13_43_bloom.png} &
        \includegraphics[width=0.17\linewidth]{../tmbe_deblooming-short-traj-latex-main/images_1/72_ordinary_CaP_40_minus_artery_50mA_5_29_bloom.png} \\[-2pt]

        \shortstack[l]{\small\textbf{without}\\\small\textbf{blooming}\\\small\textbf{artifact}} &
        \includegraphics[width=0.17\linewidth]{../tmbe_deblooming-short-traj-latex-main/images_1/03_ordinary_CaP_40_minus_artery_100mA_10_42_debloom.png} &
        \includegraphics[width=0.17\linewidth]{../tmbe_deblooming-short-traj-latex-main/images_1/32_ordinary_CaP_50_minus_artery_300mA_12_37_debloom.png} &
        \includegraphics[width=0.17\linewidth]{../tmbe_deblooming-short-traj-latex-main/images_1/35_ordinary_CaP_50_minus_artery_100mA_13_43_debloom.png} &
        \includegraphics[width=0.17\linewidth]{../tmbe_deblooming-short-traj-latex-main/images_1/72_ordinary_CaP_40_minus_artery_50mA_5_29_debloom.png} \\[-3pt]
    \end{tabular}

    \caption{Representative paired examples with and without calcium blooming artifact~\cite{shao2026deblooming}. Red circles highlight the calcified regions.}
    \label{fig:bloom_debloom_grid}
\end{figure}
U-Net~\cite{ronneberger2015unetconvolutionalnetworksbiomedical} and U-Net++~\cite{zhou2018unetnestedunetarchitecture} are widely used segmentation models and have also been adopted as encoder--decoder backbones for image restoration. In these settings, they are often trained with pixel-wise reconstruction losses such as $L_1$ and $L_2$, which promote intensity agreement but can suppress high-frequency details and produce over-smoothed outputs. This is especially limiting in CT restoration, where fine anatomical structures and sharp tissue boundaries are important. To improve perceptual sharpness, generative adversarial networks (GANs)~\cite{goodfellow2014generativeadversarialnetworks} have been explored because adversarial learning encourages generated images to match the real-image distribution beyond pixel-wise consistency. GAN-based methods have shown promising results in CT super-resolution~\cite{CT_SRGAN,Zhang2021_CTSR}, low-dose CT denoising~\cite{Yang2018_wGAN,yi2018sharpness}, domain adaptation~\cite{DIRTGAN,zhang2022generalizedoneshotdomainadaptation}, and more recently, calcium deblooming~\cite{ZHAO_sem}.

Despite these advances, existing learning-based deblooming methods may still face important limitations. First, adversarial learning can improve apparent sharpness, but GAN-based restoration may also hallucinate high-frequency edges or unstable textures that do not correspond to anatomically accurate calcium recovery. This is especially problematic in calcium deblooming, where the goal is to separate the true calcified core from the surrounding blooming spread rather than simply sharpen the plaque boundary. Even with a calcium spatial prior, the model may still preserve part of the blooming halo as calcium or suppress true plaque content. Second, many deblooming models perform restoration in a single step, requiring the network to correct calcium extent and recover the adjacent lumen simultaneously. This one-step formulation is highly ill-posed and can lead to either residual blooming or excessive calcium removal. These limitations suggest the need for a more structured restoration strategy that progressively reduces blooming while preserving calcium morphology and lumen consistency.
More recently, diffusion models have become effective tools for medical image analysis and CT restoration, including denoising~\cite{xia2022lowdosectusingdenoising,liuZeroShotdenoising}, reconstruction~\cite{Peng2024}, super-resolution~\cite{niu20243dphotoncountingct,wiedeman2024photoncountingctusingconditional}, and other restoration tasks~\cite{Graf2023,hidiff,LatentPaint,liu2025erasediffusionempoweringobject,ShiPrivacy,wang2025pathologyinformedlatentdiffusionmodel,WODZINSKI2024109129}. However, conventional diffusion models usually recover clean images from pure noise through a long reverse trajectory, which is computationally expensive and may not sufficiently constrain the output for a task-specific restoration problem. In contrast, ResShift~\cite{resshift} defines a shorter trajectory between low- and high-quality images, making it more suitable for image restoration. This is well aligned with calcium deblooming, where the objective is not unconditional generation, but progressive image-domain reduction of blooming artifacts from a corrupted CCTA input toward a debloomed target.

An effective deblooming framework should satisfy three key requirements: it should progressively reduce the artifact through a structured restoration process rather than a single direct mapping, focus the correction on calcium-associated regions instead of unnecessarily modifying the entire image, and incorporate boundary-aware modeling to better separate true calcium from the adjacent annulus. Motivated by these requirements, we propose \textbf{Short-Trajectory Artifact-Reversal (STAR) Diffusion} for CCTA calcium deblooming. Here, artifact-reversal is used in the image domain: the model learns to reduce blooming-corrupted appearance in CCTA patches and recover a debloomed target image representation.
The main contributions of this work are summarized as follows:
\begin{itemize}
    \item We formulate CCTA calcium deblooming as a short-trajectory image restoration problem and propose STAR Diffusion to progressively restore blooming-corrupted patches toward debloomed targets while preserving calcium morphology and adjacent lumen structure.

    \item We introduce an ROI-localized corruption and reverse-restoration process that uses the calcium spatial prior to restrict the trajectory to calcium-associated regions, focusing model capacity on the blooming artifact and improving inference efficiency.

    \item We introduce Soft Annular Boundary Attention (SABA), a Swin attention modification that applies a directional soft penalty from calcium-core query tokens to annular-ring key tokens, reducing boundary leakage and improving separation of true calcium from the blooming annulus.
\end{itemize}


\section{Related Works}
Prior CT calcium deblooming studies include image-domain post-processing, multi-reconstruction fusion, subtraction-based correction, and reconstruction-domain optimization. Blind deconvolution with anisotropic diffusion filtering sharpens calcium boundaries while suppressing amplified noise, but remains sensitive to blur-model estimation and sharpening artifacts~\cite{spie_conv}. Resolution-fusion methods combine smooth and sharp reconstructions to improve plaque boundaries while preserving noise texture~\cite{smooth_sharp}, but require multiple reconstructions of the same acquisition. Subtraction-based methods reduce stenosis overestimation using filtering, clustering, or registered non-contrast/contrast CT pairs~\cite{pp_subtraction,wall_calc}; however, they may be sensitive to intensity overlap, registration errors, and additional scan requirements. Reconstruction-domain methods such as model-based algebraic iterative reconstruction address blooming closer to image formation~\cite{mbai}, but require projection data and rely on a calcium/background decomposition that may be too idealized for real CCTA. Overall, conventional deblooming approaches often depend on multiple image inputs or projection data, motivating single-image learning-based methods that operate directly on reconstructed CT images while mitigating the artifact and noise trade-offs of traditional single-image post-processing.

Recent learning-based studies have begun to address calcium deblooming directly from reconstructed CCTA images. Sun et al. treated deblooming as a super-resolution problem using a fine-tuned SRGAN framework to recover high-frequency plaque boundary details~\cite{sun_ng_2022}; however, super-resolution GANs may hallucinate high-frequency structures rather than faithfully recover true plaque or lumen morphology. Zhao et al. proposed a semantic-oriented GAN that exploits calcified plaques as semantic regions through global--local fusion and semantic similarity modeling~\cite{ZHAO_sem}, while Shao et al. introduced a simulation-based paired-data framework using CatSim training data and incorporated a calcium mask as an additional input channel for U-Net++-based deblooming~\cite{shao2026deblooming}. Related anatomically guided CAC segmentation work further shows that heart and vessel localization can improve calcium-region plausibility~\cite{bioengineering12111243casc_seg}. Nevertheless, one-step pixel-loss restoration methods may suppress blooming artifacts without reliably recovering plaque morphology or the surrounding lumen. These limitations motivate short-trajectory, spatially guided, boundary-aware deblooming that preserves the calcified core while reducing the blooming-affected annulus.


\begin{figure}[t]
    \centering
    \includegraphics[width=0.95\linewidth]{../tmbe_deblooming-short-traj-latex-main/images/Slide1.PNG}
    \caption{Overview of STAR Diffusion for CCTA deblooming. The forward process shifts a debloomed target patch toward a blooming-corrupted input patch, and reverse inference progressively recovers the clean target. Refer to Sec.~\ref{sec:methods} for more details.}
    \label{fig:proposed_framework}
\end{figure}

\section{Methods}
\label{sec:methods}
As shown in Fig.~\ref{fig:proposed_framework}, the proposed method consists of STAR Diffusion, a short-trajectory framework that progressively reverses an ROI-localized corruption path from debloomed target patches to blooming-corrupted input patches. This enables efficient stepwise restoration from blooming-corrupted CCTA to debloomed CCTA. The framework also incorporates \textbf{Soft Annular Boundary Attention (SABA)}, which applies a directional soft penalty from calcium-core query tokens to annular-ring key tokens at bottleneck-adjacent decoder attention stages, helping separate true calcium from blooming artifact and supporting lumen recovery.

\subsection{Short-Trajectory Chain to Reverse CCTA Blooming}

The proposed diffusion formulation is inspired by short-trajectory restoration models such as ResShift~\cite{resshift}, but is adapted to calcium-localized CCTA blooming. In ResShift-style restoration, the forward process shifts from a high-quality image to a low-quality counterpart. Here, it shifts from a debloomed target patch to a blooming-corrupted input patch. Since blooming is localized around calcified plaques, we restrict corruption and reverse restoration to a calcium-centered ROI. Let $x_0$ denote the debloomed target patch and $x_T$ the blooming-corrupted input endpoint. Here, $T$ is the terminal degradation timestep, so $x_T$ initializes reverse inference. At each reverse step, the network predicts the clean target estimate $\hat{x}_0$, and the sampled next state $x_{t-1}$ is obtained analytically from the posterior. The calcium spatial prior $m$ localizes the blooming-prone region and guides restoration toward the calcium boundary and adjacent lumen.

\subsubsection{ROI-Localized Forward Process}

Let $m \in [0,1]^{H \times W}$ denote the calcium spatial prior. We define the spatial-prior-weighted residual between the blooming-corrupted input patch and the debloomed target patch as
\[
r_m = m \odot (x_T - x_0),
\]
where $\odot$ denotes element-wise multiplication. The forward process gradually adds this calcium-localized residual to the debloomed target patch through a Markov chain:
\[
\begin{aligned}
q(x_t \mid x_{t-1}, x_0, x_T, m)
&=
\mathcal{N}\!\left(
x_t;\,
x_{t-1}+\alpha_t r_m,\,
\kappa^2 \alpha_t I
\right),\\
&\quad t=1,2,\ldots,T .
\end{aligned}
\]
where $\alpha_t=\eta_t-\eta_{t-1}$ controls the residual-shifting step size, $\eta_0=0$, $\eta_T=1$, and $\kappa$ controls the noise variance. This formulation moves $x_0$ toward $x_T$ only inside the calcium-related ROI, preserving the rest of the patch.

The marginal distribution of this process can be written as
\[
q(x_t \mid x_0, x_T, m)
=
\mathcal{N}\!\left(
x_t;\,
x_0+\eta_t r_m,\,
\kappa^2 \eta_t I
\right).
\]
Thus, as $t$ increases, the patch follows a short residual trajectory from the debloomed target toward the blooming-corrupted input endpoint. The schedule is controlled by $\{\eta_t\}_{t=1}^{T}$ and $\kappa$, with $\eta_t$ increasing from $0$ to $1$.



\subsubsection{Reverse Deblooming Process}

The reverse process approximates the inverse of the ROI-localized corruption trajectory. At inference, the chain is initialized from the observed blooming-corrupted input patch $x_T \equiv x_{\mathrm{bloom}}$. Following the ResShift-style sampling procedure~\cite{resshift}, the network predicts the clean debloomed target at each step and samples the next reverse state from the corresponding Gaussian posterior. The detailed update is given in Supplementary Sec.~\ref{supp:reverse_sampling}.

\[
\hat{x}_0
=
f_{\theta}(x_t, m, t).
\]

This keeps the learning target fixed across timesteps and leaves the timestep-specific transition to the analytic reverse kernel.

The model is optimized using a patch restoration objective together with ROI- and boundary-focused losses. The whole-patch restoration loss is defined using an $L_1$ objective:

\[
\mathcal{L}_{\mathrm{patch}}
=
\left\|
\hat{x}_0 - x_0
\right\|_1 .
\]

To emphasize the calcium-localized blooming region, we additionally use an ROI-focused reconstruction loss based on the same calcium spatial prior $m$:

\[
\mathcal{L}_{\mathrm{ROI}}
=
\left\|
m \odot
(\hat{x}_0-x_0)
\right\|_1 ,
\]

Here, $m$ is the dilated calcium spatial prior used for conditioning and for ROI-localized supervision.

To explicitly supervise the bloom halo and plaque-lumen boundary, we derive an annular spatial-prior band around the calcium region:

\[
m_{\mathrm{ann}}
=
\mathrm{Dilate}(m,1)
-
\mathrm{Erode}(m,1).
\]
This band captures the immediate calcium-lumen neighborhood, where residual blooming and overcorrection are most likely to appear.

Instead of applying a direct reconstruction loss in this annular region, we use an annulus-focused Learned Perceptual Image Patch Similarity (LPIPS) perceptual loss~\cite{zhang2018unreasonableeffectivenessdeepfeatures}:

\[
\begin{aligned}
\mathcal{L}_{\mathrm{ann\text{-}perc}}
=
\mathrm{LPIPS}\Big(
&m_{\mathrm{ann}}\odot \hat{x}_0 \\
&+ (1-m_{\mathrm{ann}})\odot x_0,\ x_0
\Big).
\end{aligned}
\]

Finally, an edge-consistency loss is used to encourage sharp plaque and lumen boundaries. We define a narrow edge band from the spatial prior as

\[
m_{\mathrm{edge}}
=
\mathrm{Dilate}\!\left(\mathbf{1}\{|\nabla m|>0\},1\right),
\]

and use it to localize the gradient penalty:

\[
\mathcal{L}_{\mathrm{edge}}
=
\left\|
m_{\mathrm{edge}} \odot
\left(
\nabla \hat{x}_0
-
\nabla x_0
\right)
\right\|_1 ,
\]

where $m_{\mathrm{edge}}$ differs from $m_{\mathrm{ann}}$ by focusing only on the immediate boundary pixels used for gradient consistency, while $m_{\mathrm{ann}}$ covers the broader annular halo used for perceptual supervision. The final optimization objective is

\[
\begin{aligned}
\min_{\theta}\ \mathcal{L}_{\mathrm{total}}
=&\
\lambda_{\mathrm{patch}}\mathcal{L}_{\mathrm{patch}}
+
\lambda_{\mathrm{ROI}}\mathcal{L}_{\mathrm{ROI}} \\
&+
\lambda_{\mathrm{perc}}\mathcal{L}_{\mathrm{ann\text{-}perc}}
+
\lambda_{\mathrm{edge}}\mathcal{L}_{\mathrm{edge}} .
\end{aligned}
\]

By learning this short, ROI-localized reverse chain with spatially guided and boundary-focused supervision, the model performs deblooming as a gradual inversion of the calcium blooming process rather than a single-step correction. This design is better aligned with the physical nature of blooming, where high-density calcium spreads into neighboring lumen regions, and allows the network to progressively reduce the blooming halo while preserving the true calcium core and recovering the obscured lumen.


\begin{figure}[t]
    \centering
    \includegraphics[width=0.95\linewidth]{../tmbe_deblooming-short-traj-latex-main/images/Slide2.PNG}
    \caption{Illustration of the annular boundary-guided Swin attention mechanism used in SABA. The resized calcium spatial prior defines the core and annular boundary regions, and the annulus is used to penalize attention from calcium-core query tokens to annular-ring key tokens. Refer to Sec.~\ref{sec:saba} for more details.}
    \label{fig:saba_attention}
\end{figure}

\subsection{Soft Annular Boundary Attention}
\label{sec:saba}

Although the short-trajectory diffusion chain provides stepwise restoration, accurate deblooming still requires control of feature interactions near the calcium-lumen boundary. As illustrated in Fig.~\ref{fig:saba_attention}, we form an annular boundary zone by resizing the spatial prior to the deep Swin feature resolution. Standard Swin window attention lets calcium-core tokens attend to nearby annular tokens, which can mix true calcium with blooming artifact. To reduce this boundary leakage, we introduce \textbf{Soft Annular Boundary Attention (SABA)}, a lightweight Swin attention modification that applies a directional soft penalty from calcium-core query tokens to annular-ring key tokens.

At feature level $l$, let $m$ denote the calcium spatial prior and $S_{\mathrm{dilated}}^{(l)}$ its resized version:
\[
S_{\mathrm{dilated}}^{(l)}
=
\mathrm{Resize}(m,H_l,W_l).
\]
Since the spatial prior is already dilated, we estimate the calcium core by applying a one-step local erosion operation:
\[
S_{\mathrm{core}}^{(l)}
=
\mathrm{Erode}_{k}
\left(
S_{\mathrm{dilated}}^{(l)}
\right),
\]
where $k=1$ in this paper. The annular boundary ring is then obtained by subtracting the estimated core from the spatial-prior map:
\[
A_{\mathrm{ring}}^{(l)}
=
S_{\mathrm{dilated}}^{(l)}
-
S_{\mathrm{core}}^{(l)}.
\]
Thus, $S_{\mathrm{core}}^{(l)}$ identifies calcium-core tokens, while $A_{\mathrm{ring}}^{(l)}$ identifies the immediate boundary region around the calcium core.

Given the standard Swin window attention formulation,
\[
A
=
\mathrm{Softmax}
\left(
\frac{QK^{T}}{\sqrt{d}}
+
R
+
M
\right),
\]
where $R$ is the learned relative-position bias and $M$ is the shifted-window attention mask. SABA adds an annular penalty before softmax. For query token $i$ and key token $j$ in the same window, the annular bias is
\[
B_{ij}
=
-\lambda S_i^{(l)} A_{\mathrm{ring},j}^{(l)},
\]
where $S_i^{(l)}=1$ indicates that query token $i$ belongs to the calcium core, $A_{\mathrm{ring},j}^{(l)}=1$ indicates that key token $j$ belongs to the annular ring, and $\lambda>0$ controls the strength of boundary suppression.

The modified Swin attention becomes
\[
A_{ij}
=
\mathrm{Softmax}_{j}
\left(
\frac{q_i^{T}k_j}{\sqrt{d}}
+
R_{ij}
+
M_{ij}
-
\lambda S_i^{(l)} A_{\mathrm{ring},j}^{(l)}
\right).
\]
This penalty is directional: it suppresses calcium-core queries to annular-ring keys without blocking the reverse direction. Because the penalty is finite, SABA reduces artifact-contaminated links by a factor proportional to $\exp(-\lambda)$ while preserving useful context.

SABA is applied on the decoder side near the bottleneck, while the encoder keeps standard Swin attention to preserve early feature extraction. In our architecture, it is used at the bottleneck and the adjacent decoder $8\times8$ and $16\times16$ stages, where boundary reconstruction is most sensitive to annulus-to-core leakage and therefore benefits most from the penalty.
\begin{figure}[t]
    \centering
    \setlength{\tabcolsep}{1.2pt}
    \renewcommand{\arraystretch}{1.05}
    \scriptsize

    % ---------------- Case 1 ----------------
    \begin{tabular}{ccccccc}
        \textbf{Input} &
        \shortstack{\textbf{Morphological}\\\textbf{Shrinking}} &
        \shortstack{\textbf{Wiener}\\\textbf{Deconvolution}} &
        \textbf{SRGAN} &
        \textbf{SRGAN-no-style} &
        \textbf{SAGAN} &
        \textbf{SOGAN} \\

        
\includegraphics[width=0.12\linewidth]{../tmbe_deblooming-short-traj-latex-main/images_sim/03_CaP_50_minus_artery_100mA_2_slice32_input.png} &
        \includegraphics[width=0.12\linewidth]{../tmbe_deblooming-short-traj-latex-main/images_sim/03_CaP_50_minus_artery_100mA_2_slice32_HU_Shrink.png} &
        \includegraphics[width=0.12\linewidth]{../tmbe_deblooming-short-traj-latex-main/images_sim/03_CaP_50_minus_artery_100mA_2_slice32_Wiener.png} &
        
\includegraphics[width=0.12\linewidth]{../tmbe_deblooming-short-traj-latex-main/images_sim/03_CaP_50_minus_artery_100mA_2_slice32_pred_output_srgan.png} &
        
\includegraphics[width=0.12\linewidth]{../tmbe_deblooming-short-traj-latex-main/images_sim/03_CaP_50_minus_artery_100mA_2_slice32_pred_output_srgan-no_style.png} &
        
\includegraphics[width=0.12\linewidth]{../tmbe_deblooming-short-traj-latex-main/images_sim/03_CaP_50_minus_artery_100mA_2_slice32_pred_output_sagan.png} &
        
\includegraphics[width=0.12\linewidth]{../tmbe_deblooming-short-traj-latex-main/images_sim/03_CaP_50_minus_artery_100mA_2_slice32_pred_output_sogan.png} \\

        \textbf{Spatial Prior} &
        \textbf{U-Net++} &
        \textbf{U-Net++ Sharp} &
        \textbf{LDM-1000} &
        \textbf{LDM-10} &
        \textbf{Proposed} &
        \textbf{GT} \\

        
\includegraphics[width=0.12\linewidth]{../tmbe_deblooming-short-traj-latex-main/images_sim/03_CaP_50_minus_artery_100mA_2_slice32_mask.png} &
        
\includegraphics[width=0.12\linewidth]{../tmbe_deblooming-short-traj-latex-main/images_sim/03_CaP_50_minus_artery_100mA_2_slice32_pred_output_unetpp_amp.png} &
        
\includegraphics[width=0.12\linewidth]{../tmbe_deblooming-short-traj-latex-main/images_sim/03_CaP_50_minus_artery_100mA_2_slice32_pred_output_unetpp_sharp_patient.png} &
        \includegraphics[width=0.12\linewidth]{../tmbe_deblooming-short-traj-latex-main/images_sim/03_CaP_50_minus_artery_100mA_2_slice32_pred_ldm_output_resshift.png} &
        \includegraphics[width=0.12\linewidth]{../tmbe_deblooming-short-traj-latex-main/images_sim/03_CaP_50_minus_artery_100mA_2_slice32_pred_ldm10_output_resshift.png} &
        \includegraphics[width=0.12\linewidth]{../tmbe_deblooming-short-traj-latex-main/images_sim/03_CaP_50_minus_artery_100mA_2_slice32_pred_output_resshift.png} &
        
\includegraphics[width=0.12\linewidth]{../tmbe_deblooming-short-traj-latex-main/images_sim/03_CaP_50_minus_artery_100mA_2_slice32_gt.png}
    \end{tabular}

    \vspace{0.8em}

    % ---------------- Case 2 ----------------
    \begin{tabular}{ccccccc}
        \textbf{Input} &
        \shortstack{\textbf{Morphological}\\\textbf{Shrinking}} &
        \shortstack{\textbf{Wiener}\\\textbf{Deconvolution}} &
        \textbf{SRGAN} &
        \textbf{SRGAN-no-style} &
        \textbf{SAGAN} &
        \textbf{SOGAN} \\

        
\includegraphics[width=0.12\linewidth]{../tmbe_deblooming-short-traj-latex-main/images_sim/03_CaP_50_minus_artery_100mA_10_slice32_input.png} &
        \includegraphics[width=0.12\linewidth]{../tmbe_deblooming-short-traj-latex-main/images_sim/03_CaP_50_minus_artery_100mA_10_slice32_HU_Shrink.png} &
        \includegraphics[width=0.12\linewidth]{../tmbe_deblooming-short-traj-latex-main/images_sim/03_CaP_50_minus_artery_100mA_10_slice32_Wiener.png} &
        
\includegraphics[width=0.12\linewidth]{../tmbe_deblooming-short-traj-latex-main/images_sim/03_CaP_50_minus_artery_100mA_10_slice32_pred_output_srgan.png} &
        
\includegraphics[width=0.12\linewidth]{../tmbe_deblooming-short-traj-latex-main/images_sim/03_CaP_50_minus_artery_100mA_10_slice32_pred_output_srgan-no_style.png} &
        
\includegraphics[width=0.12\linewidth]{../tmbe_deblooming-short-traj-latex-main/images_sim/03_CaP_50_minus_artery_100mA_10_slice32_pred_output_sagan.png} &
        
\includegraphics[width=0.12\linewidth]{../tmbe_deblooming-short-traj-latex-main/images_sim/03_CaP_50_minus_artery_100mA_10_slice32_pred_output_sogan.png} \\

        \textbf{Spatial Prior} &
        \textbf{U-Net++} &
        \textbf{U-Net++ Sharp} &
        \textbf{LDM-1000} &
        \textbf{LDM-10} &
        \textbf{Proposed} &
        \textbf{GT} \\

        
\includegraphics[width=0.12\linewidth]{../tmbe_deblooming-short-traj-latex-main/images_sim/03_CaP_50_minus_artery_100mA_10_slice32_mask.png} &
        
\includegraphics[width=0.12\linewidth]{../tmbe_deblooming-short-traj-latex-main/images_sim/03_CaP_50_minus_artery_100mA_10_slice32_pred_output_unetpp_amp.png} &
        
\includegraphics[width=0.12\linewidth]{../tmbe_deblooming-short-traj-latex-main/images_sim/03_CaP_50_minus_artery_100mA_10_slice32_pred_output_unetpp_sharp_patient.png} &
        \includegraphics[width=0.12\linewidth]{../tmbe_deblooming-short-traj-latex-main/images_sim/03_CaP_50_minus_artery_100mA_10_slice32_pred_ldm_output_resshift.png} &
        \includegraphics[width=0.12\linewidth]{../tmbe_deblooming-short-traj-latex-main/images_sim/03_CaP_50_minus_artery_100mA_10_slice32_pred_ldm10_output_resshift.png} &
        \includegraphics[width=0.12\linewidth]{../tmbe_deblooming-short-traj-latex-main/images_sim/03_CaP_50_minus_artery_100mA_10_slice32_pred_output_resshift.png} &
        
\includegraphics[width=0.12\linewidth]{../tmbe_deblooming-short-traj-latex-main/images_sim/03_CaP_50_minus_artery_100mA_10_slice32_gt.png}
    \end{tabular}

    \caption{Qualitative comparison on synthetic CCTA patches. The first row shows the input and baseline methods; the second row shows the spatial prior, latent diffusion baselines, the proposed method, and ground truth. All patches are displayed using WW/WL = 1300/400 HU.}
    \label{fig:synthetic_qualitative}
\end{figure}


\begin{table}[t]
\centering
\caption{Quantitative comparison of calcium deblooming performance. Dice is computed on the calcium core using $>850$ HU, and Core RMSE and Annulus RMSE are computed in the core and annulus, respectively.}
\label{tab:deblooming_results}
\scriptsize
\setlength{\tabcolsep}{8pt}
\begin{tabular}{lccc}
\hline
\textbf{Method} &
\textbf{Dice $\uparrow$} &
\begin{tabular}[c]{@{}c@{}}\textbf{Core}\\\textbf{RMSE $\downarrow$}\end{tabular} &
\begin{tabular}[c]{@{}c@{}}\textbf{Annulus}\\\textbf{RMSE $\downarrow$}\end{tabular} \\
\hline
Blooming-corrupted input & $0.3842 \pm 0.3366$ & $341.4590 \pm 34.8052$ & $176.1679 \pm 14.7582$ \\
Morphological Shrinking & $0.3842 \pm 0.2590$ & $746.4000 \pm 103.8000$ & $508.0000 \pm 99.2000$ \\
Wiener Deconvolution & $0.7082 \pm 0.1381$ & $225.5000 \pm 51.6000$ & $231.0000 \pm 60.2000$ \\
SRGAN & $0.7405 \pm 0.1046$ & $215.3337 \pm 46.8714$ & $212.0082 \pm 41.3469$ \\
SRGAN-no-style & $0.6766 \pm 0.1446$ & $245.6082 \pm 55.6651$ & $215.4069 \pm 45.3816$ \\
SAGAN & $0.6792 \pm 0.1372$ & $249.6572 \pm 56.2867$ & $225.5519 \pm 45.2428$ \\
SOGAN & $\underline{0.8922} \pm 0.0587$ & $\textbf{131.4178} \pm 19.5148$ & $155.5562 \pm 32.3092$ \\
U-Net++ & $0.8870 \pm 0.0284$ & $187.1475 \pm 27.7500$ & $\underline{144.0696} \pm 23.3815$ \\
U-Net++ Sharp & $0.8171 \pm 0.0412$ & $153.7231 \pm 17.5009$ & $235.9129 \pm 35.7622$ \\
LDM-1000 & $0.7205 \pm 0.1807$ & $278.3594 \pm 33.2693$ & $168.9410 \pm 16.1816$ \\
LDM-10 & $0.7071 \pm 0.1342$ & $311.0971 \pm 32.5633$ & $176.0186 \pm 22.4940$ \\
Proposed & $\textbf{0.9250} \pm 0.0781$ & $\underline{134.4998} \pm 32.2103$ & $\textbf{90.2365} \pm 9.7464$ \\
\hline
\end{tabular}
\end{table}

\begin{figure}[t]
    \centering
    \setlength{\tabcolsep}{1.5pt}
    \renewcommand{\arraystretch}{1.05}

    % ---------------- Patient 1 ----------------
    \begin{tabular}{ccccccc}
        \textbf{Input} &
        \shortstack{\textbf{Morphological}\\\textbf{Shrinking}} &
        \shortstack{\textbf{Wiener}\\\textbf{Deconvolution}} &
        \textbf{SRGAN} &
        \textbf{SRGAN-no-style} &
        \textbf{SAGAN} &
        \textbf{SOGAN} \\

        
\includegraphics[width=0.12\linewidth]{../tmbe_deblooming-short-traj-latex-main/images_real/e29056_187_region_1_input_raw_sagan.png} &
        \includegraphics[width=0.12\linewidth]{../tmbe_deblooming-short-traj-latex-main/images_real/e29056_187_region_1_HU_Shrink.png} &
        \includegraphics[width=0.12\linewidth]{../tmbe_deblooming-short-traj-latex-main/images_real/e29056_187_region_1_Wiener_Unsharp.png} &
        
\includegraphics[width=0.12\linewidth]{../tmbe_deblooming-short-traj-latex-main/images_real/e29056_187_region_1_output_raw_srgan.png} &
        
\includegraphics[width=0.12\linewidth]{../tmbe_deblooming-short-traj-latex-main/images_real/e29056_187_region_1_output_raw_srgan_no_style.png} &
        
\includegraphics[width=0.12\linewidth]{../tmbe_deblooming-short-traj-latex-main/images_real/e29056_187_region_1_output_raw_sagan.png} &
        \includegraphics[width=0.12\linewidth]{../tmbe_deblooming-short-traj-latex-main/images_real/e29056_187_region_1_output_raw_sogan_mask_undilate.png} \\

        \textbf{Spatial Prior} &
        \textbf{U-Net++} &
        \textbf{U-Net++ Sharp} &
        \textbf{LDM-1000} &
        \textbf{LDM-10} &
        \textbf{Proposed} &
        {} \\

        
\includegraphics[width=0.12\linewidth]{../tmbe_deblooming-short-traj-latex-main/images_real/e29056_187_region_1_mask_raw.png} &
        
\includegraphics[width=0.12\linewidth]{../tmbe_deblooming-short-traj-latex-main/images_real/e29056_187_region_1_output_unetpp_amp.png} &
        
\includegraphics[width=0.12\linewidth]{../tmbe_deblooming-short-traj-latex-main/images_real/e29056_187_region_1_output_unetpp_sharp_patient.png} &
        \includegraphics[width=0.12\linewidth]{../tmbe_deblooming-short-traj-latex-main/images_real/e29056_187_region_1_pred_ldm_output_resshift.png} &
        \includegraphics[width=0.12\linewidth]{../tmbe_deblooming-short-traj-latex-main/images_real/e29056_187_region_1_pred_ldm10_output_resshift.png} &
        \includegraphics[width=0.12\linewidth]{../tmbe_deblooming-short-traj-latex-main/images_real/e29056_187_region_1_output_resshift.png} &
        {}
    \end{tabular}

    \vspace{0.7em}

    % ---------------- Patient 2 ----------------
    \begin{tabular}{ccccccc}
        \textbf{Input} &
        \shortstack{\textbf{Morphological}\\\textbf{Shrinking}} &
        \shortstack{\textbf{Wiener}\\\textbf{Deconvolution}} &
        \textbf{SRGAN} &
        \textbf{SRGAN-no-style} &
        \textbf{SAGAN} &
        \textbf{SOGAN} \\

        
\includegraphics[width=0.12\linewidth]{../tmbe_deblooming-short-traj-latex-main/images_real/e35127_147_region_1_input_raw_sagan.png} &
        \includegraphics[width=0.12\linewidth]{../tmbe_deblooming-short-traj-latex-main/images_real/e35127_147_region_1_HU_Shrink.png} &
        \includegraphics[width=0.12\linewidth]{../tmbe_deblooming-short-traj-latex-main/images_real/e35127_147_region_1_Wiener_Unsharp.png} &
        
\includegraphics[width=0.12\linewidth]{../tmbe_deblooming-short-traj-latex-main/images_real/e35127_147_region_1_output_raw_srgan.png} &
        
\includegraphics[width=0.12\linewidth]{../tmbe_deblooming-short-traj-latex-main/images_real/e35127_147_region_1_output_raw_srgan_no_style.png} &
        
\includegraphics[width=0.12\linewidth]{../tmbe_deblooming-short-traj-latex-main/images_real/e35127_147_region_1_output_raw_sagan.png} &
        
\includegraphics[width=0.12\linewidth]{../tmbe_deblooming-short-traj-latex-main/images_real/e35127_147_region_1_output_sogan_mask_undilate.png} \\

        \textbf{Spatial Prior} &
        \textbf{U-Net++} &
        \textbf{U-Net++ Sharp} &
        \textbf{LDM-1000} &
        \textbf{LDM-10} &
        \textbf{Proposed} &
        {} \\

        \includegraphics[width=0.12\linewidth]{../tmbe_deblooming-short-traj-latex-main/images_real/e35127_147_region_1_mask_raw.png} &
        
\includegraphics[width=0.12\linewidth]{../tmbe_deblooming-short-traj-latex-main/images_real/e35127_147_region_1_output_unetpp_amp.png} &
        
\includegraphics[width=0.12\linewidth]{../tmbe_deblooming-short-traj-latex-main/images_real/e35127_147_region_1_output_unetpp_sharp_patient.png} &
        \includegraphics[width=0.12\linewidth]{../tmbe_deblooming-short-traj-latex-main/images_real/e35127_147_region_1_pred_ldm_output_resshift.png} &
        \includegraphics[width=0.12\linewidth]{../tmbe_deblooming-short-traj-latex-main/images_real/e35127_147_region_1_pred_ldm10_output_resshift.png} &
        \includegraphics[width=0.12\linewidth]{../tmbe_deblooming-short-traj-latex-main/images_real/e35127_147_region_1_output_resshift.png} &
        {}
    \end{tabular}

    \caption{Qualitative comparison on clinical CCTA patches. The first row shows the input and baseline methods; the second row shows the spatial prior, latent diffusion baselines, and the proposed method. All patches are displayed using WW/WL = 1300/400 HU.}
    \label{fig:real_patient_qualitative}
\end{figure}


\section{Experiments}

\subsection{Experimental Setup}

\noindent\textbf{Datasets.}

We evaluate the proposed method using the synthetic CCTA deblooming dataset from Shao et al.~\cite{shao2026deblooming}, which provides paired blooming-corrupted and debloomed $128 \times 128 \times 64$ CCTA volumes with calcium spatial priors. The dataset contains 3{,}463 MAT files from 12 source subjects, with multiple anatomical backgrounds, three calcium concentrations, and four exposure/noise levels. We use a subject-level split with no overlap between training, validation, and testing: 8 subjects for training, 2 for validation, and 2 for testing, corresponding to 2{,}402, 365, and 696 MAT files, respectively.
The model is trained slice-wise using 20{,}000 random $128 \times 128$ patches from the training subjects and 2{,}500 patches from the validation subjects. For final testing, we perform volume-level evaluation on 232 held-out 3D cases selected from the 696 test volumes. Each selected case represents a distinct anatomical background, calcium concentration, and exposure/noise setting, avoiding repeated slices or duplicate test conditions. Inference is applied slice-by-slice and reassembled into 3D volumes for evaluation. Quantitative evaluation is performed only on synthetic data because paired references are available; clinical CCTA data from the same authors are used only for qualitative assessment.

\noindent \textbf{Baselines.}
We compare the proposed method with classical, single-step learning-based, and diffusion-based baselines. The classical baselines include morphological shrinking and Wiener deconvolution. The learning-based baselines include U-Net++~\cite{zhou2018unetnestedunetarchitecture}, U-Net++ Sharp, SRGAN~\cite{CT_SRGAN}, SRGAN-no-style, SAGAN~\cite{yi2018sharpness}, and SOGAN~\cite{ZHAO_sem}. SRGAN-no-style removes the VGG Gram-matrix style loss from SRGAN, while U-Net++ Sharp adds the Sobel edge-magnitude sharpness loss used in SAGAN.

To compare against a long-trajectory diffusion baseline, we also evaluate a conditional latent diffusion model (LDM)~\cite{rombach2022highresolution}. The LDM generates debloomed images conditioned on the blooming-corrupted input and spatial prior. We use latent diffusion because it is more efficient than pixel-space diffusion and can reduce dense pixel-level learning burden for localized restoration~\cite{ju2024brushnet,LatentPaint}. We report two sampling variants from the same trained model: LDM-1000 with standard 1000-step DDPM sampling and LDM-10 with accelerated 10-step DDIM sampling.

\noindent \textbf{Evaluation Metrics.}

The main goal of calcium deblooming is to recover calcified plaque morphology while improving coronary lumen visibility. We evaluate performance using three volumetric metrics: Dice for calcium-core agreement, Core RMSE for HU fidelity within the calcified plaque, and Annulus RMSE for residual blooming or overcorrection near the plaque-lumen interface. 
Prior 3D-printed coronary plaque studies report calcified plaques around 800--850 HU with measured attenuation well above that level~\cite{biom11091307}, consistent with the dataset used by Shao et al.~\cite{shao2026deblooming}; we therefore define the calcium core using an $>850$ HU threshold, where blooming is most pronounced. The annular boundary region is defined as $550<\mathrm{HU}\leq850$. 
For each method, metrics are computed once per selected 3D test volume after slice-wise inference and volumetric reassembly, yielding 232 volume-level values for statistical comparison.

\noindent \textbf{Implementation Details.}
For the GAN- and U-Net-based baselines~\cite{CT_SRGAN,yi2018sharpness,ZHAO_sem,shao2026deblooming}, we follow the original architectures and losses and retrain all models on the same subject-level split. Models are optimized with Adam using an initial learning rate of $10^{-5}$ and exponential decay of 0.99 per epoch. Training is run for 200 epochs, and the checkpoint with the best validation loss is reported. For spatial-prior-guided variants, the blooming-corrupted CCTA patch is concatenated with the calcium spatial prior. No additional geometric or intensity augmentation is used.

For the proposed STAR Diffusion method, we use a pixel-space spatial-prior-conditioned Swin U-Net diffusion model. The blooming-corrupted CCTA patch and calcium spatial prior are provided as conditional inputs at each denoising step, and ROI-localized corruption restricts the diffusion process to calcium-centered blooming regions. The model is trained for 30{,}000 iterations with a 10-step diffusion trajectory, Adam optimization, initial learning rate $5\times10^{-5}$, cosine decay to $2\times10^{-5}$, 5{,}000 warm-up iterations, batch size 24, and EMA rate 0.999. The reported checkpoint is selected using the same held-out validation subjects.

For the LDM baseline, we train a conditional latent diffusion model to generate debloomed CCTA targets from noise, conditioned on the blooming-corrupted input and spatial prior. The same trained LDM model is evaluated using the two sampling settings described above: LDM-1000 and LDM-10. Runtime experiments are performed on an NVIDIA RTX A5000 system using PyTorch 2.1.1+cu121 and cuDNN 8902.

\subsection{Results}

\noindent\textbf{Quantitative Evaluation}
For quantitative evaluation, we use the simulated subset of the dataset from Shao et al.~\cite{shao2026deblooming}, which provides paired blooming-corrupted inputs, ground-truth debloomed targets, and calcium-region annotations. All methods are evaluated against the same ground-truth target using Dice score for calcium-region agreement, Core RMSE within the calcium core, and Annulus RMSE within the surrounding annular region. The original blooming-corrupted input provides a baseline measure of artifact severity prior to correction, showing poor calcium-region agreement with a Dice score of 0.3842 and large intensity errors in both evaluated regions. The remaining rows in Table~\ref{tab:deblooming_results} compare the deblooming methods, including the proposed STAR Diffusion approach.

At significance level 0.05, using an approximate Welch's two-sample $t$-test on the reported summary statistics, the proposed method was significantly better than SOGAN in Dice ($p=2.39\times10^{-7}$) and Annulus RMSE ($p=4.39\times10^{-87}$), while Core RMSE was not significantly different from SOGAN ($p\approx0.21$). 
The main observations from Table~\ref{tab:deblooming_results} are as follows. i) The classical baselines are insufficient for this task: morphological shrinking substantially degrades both calcium morphology and HU fidelity, while Wiener deconvolution improves over simple shrinking but remains behind the learning-based methods. ii) Methods without spatial-prior guidance, including SRGAN, SRGAN-no-style, and SAGAN, show weaker deblooming performance; compared with the strongest method in this no-prior group, the proposed method improves Dice by at least 24.9\%, reduces Core RMSE by at least 37.5\%, and reduces Annulus RMSE by at least 57.4\%. iii) Adding spatial guidance improves performance, as seen from SOGAN and U-Net++, which are stronger than the unconstrained GAN baselines because they better preserve local calcium structure. iv) The LDM variants lag behind the leading spatially guided methods, suggesting that long-trajectory generation from noise is less effective for this localized restoration problem. Overall, the proposed method performs best because it combines spatial prior conditioning with a short-trajectory restoration process, allowing better calcium-core preservation and stronger suppression of the blooming halo.

\noindent\textbf{Qualitative Evaluation}
As shown in Fig.~\ref{fig:synthetic_qualitative}, the classical baselines are insufficient for calcium deblooming: morphological shrinking overcorrects the plaque and introduces discontinuities, while Wiener deconvolution leaves residual blooming. Among single-step learning baselines, spatial-prior-guided methods such as SOGAN and U-Net++ better preserve plaque shape, whereas SRGAN, SRGAN-no-style, and SAGAN mainly improve sharpness without reliably matching ground-truth morphology. The LDM variants are also unreliable for this localized restoration task; LDM-10 shows strong shape mismatch, and LDM-1000 improves over LDM-10 but still leaves the calcium region enlarged relative to the ground truth. Overall, the proposed method gives the closest synthetic result, preserving calcium volume and edge fidelity while reducing the surrounding blooming halo.

For clinical CCTA patches, where paired ground truth is unavailable, Fig.~\ref{fig:real_patient_qualitative} shows a similar ranking under more realistic conditions, with the same failure modes becoming more pronounced. The classical and LDM baselines show visible overcorrection, residual blooming, or shape mismatch, consistent with the synthetic examples. Among spatial-prior-guided methods, U-Net++ and SOGAN remain more plausible than the non-prior GAN variants, but U-Net++ can undercorrect the blooming region and SOGAN may allow boundary bleeding near the calcium edge. In comparison, the proposed method provides the best qualitative balance, suppressing the blooming halo while avoiding excessive erosion or distortion of the calcium core.

\begin{figure}[t]
    \centering
    \setlength{\tabcolsep}{2pt}
    \renewcommand{\arraystretch}{1.05}

      {\scriptsize
      \setlength{\tabcolsep}{2pt}
      \begin{tabular}{ccccc}
          \parbox[c]{0.18\linewidth}{\centering \textbf{Input}\\[2.0em]\includegraphics[width=\linewidth]{../tmbe_deblooming-short-traj-latex-main/images_mask/patient_24_100mA_2_slice_032_input_processed.png}} &
          \parbox[c]{0.18\linewidth}{\centering \textbf{Spatial Prior}\\[2.0em]
\includegraphics[width=\linewidth]{../tmbe_deblooming-short-traj-latex-main/images_mask/patient_24_100mA_2_slice_032_mask_dilate.png}} &
          \parbox[c]{0.18\linewidth}{\centering \textbf{Ground Truth}\\[2.0em]\includegraphics[width=\linewidth]{../tmbe_deblooming-short-traj-latex-main/images_mask/patient_24_100mA_2_slice_032_gt.png}} &
          \parbox[c]{0.18\linewidth}{\centering \textbf{With spatial prior}\\\includegraphics[width=\linewidth]{../tmbe_deblooming-short-traj-latex-main/images_mask/patient_24_100mA_2_slice_032_pred.png}} &
          \parbox[c]{0.18\linewidth}{\centering \textbf{Without spatial prior}\\\includegraphics[width=\linewidth]{../tmbe_deblooming-short-traj-latex-main/images_mask/patient_24_100mA_2_slice_032_pred_no_mask.png}}
      \end{tabular}}

    \caption{Spatial-prior ablation (displayed with WW/WL=1300/400)}
    \label{fig:mask_ablation}
\end{figure}

\begin{figure}[t]
    \centering
    \setlength{\tabcolsep}{3pt}
    \renewcommand{\arraystretch}{1.05}
    \begin{tabular}{cccc}
        \textbf{Input} & \textbf{\shortstack{Proposed\\version}} & \textbf{\shortstack{Baseline version\\with Swin only}} & \textbf{\shortstack{Ground\\Truth}} \\

        % -------- Synthetic example --------
        
\includegraphics[width=0.10\linewidth]{../tmbe_deblooming-short-traj-latex-main/images_saba/03_CaP_50_minus_artery_100mA_10_slice32_input.png} &
        \includegraphics[width=0.10\linewidth]{../tmbe_deblooming-short-traj-latex-main/images_saba/patient_03_100mA_10_slice_032_pred.png} &
        \includegraphics[width=0.10\linewidth]{../tmbe_deblooming-short-traj-latex-main/images_saba/patient_03_100mA_10_slice_032_pred_swin.png} &
        
\includegraphics[width=0.10\linewidth]{../tmbe_deblooming-short-traj-latex-main/images_saba/03_CaP_50_minus_artery_100mA_10_slice32_gt.png} \\

        % -------- Real example --------
        
\includegraphics[width=0.10\linewidth]{../tmbe_deblooming-short-traj-latex-main/images_saba/e29056_187_region_1_input_raw_sagan.png} &
        \includegraphics[width=0.10\linewidth]{../tmbe_deblooming-short-traj-latex-main/images_saba/e29056_187_region_1_output_resshift.png} &
        \includegraphics[width=0.10\linewidth]{../tmbe_deblooming-short-traj-latex-main/images_saba/e29056_187_region_1_pred_swin.png} &
        {}
    \end{tabular}

    \caption{SABA ablation on synthetic (top) and clinical (bottom) examples.(displayed with WW/WL=1300/400)}
    \label{fig:saba_ablation}
\end{figure}


\begin{table}[t]
\centering
\caption{Inference-time comparison of diffusion-based methods. Runtime is reported as average seconds per $128\times128\times64$ volume.}
\label{tab:runtime_comparison}
\begin{tabular}{lccc}
\hline
\textbf{Method} & \textbf{LDM-1000} & \textbf{LDM-10} & \textbf{Proposed} \\
\hline
\textbf{Time / Volume (s) $\downarrow$} & 2778.66 & 30.09 & 18.55 \\
\hline
\end{tabular}
\end{table}

\noindent\textbf{Inference-Time Efficiency.}

We compare runtime against the latent diffusion baselines because the proposed method uses a short-trajectory diffusion process. Runtime is measured per reconstructed $128\times128\times64$ volume. The proposed method requires 18.55 s, compared with 30.09 s for LDM-10 and 2778.66 s for LDM-1000. Although LDM-10 also uses 10 sampling steps, it generates from noise in latent space, whereas the proposed method performs pixel-space restoration through a short Markov trajectory from the corrupted input. 
As a result, the proposed method is approximately $1.6\times$ faster than LDM-10 and nearly $150\times$ faster than LDM-1000, while maintaining stronger deblooming performance. 
Timing was measured with \texttt{time.perf\_counter()} over the full-volume pipeline at the configured batch size, without explicit warm-up or \texttt{cudnn.benchmark}.

\subsection{Ablation Study}

We perform two ablation studies to test the main design choices in the proposed framework: (i) removing the spatial prior to evaluate whether explicit ROI guidance is needed for localized restoration, and (ii) replacing SABA with standard Swin attention to evaluate whether the annular boundary penalty improves calcium-lumen separation. These experiments assess the separate roles of spatial localization and boundary-aware feature control in preserving the calcium core while reducing blooming artifacts.

\begin{figure}[t]
    \centering
    \setlength{\tabcolsep}{2pt}
    \renewcommand{\arraystretch}{1.05}
    \scriptsize

    \begin{minipage}[t]{0.49\linewidth}
        \centering
        \textbf{Case 1}\\[0.25em]
        \begin{tabular}{cccc}
            \textbf{Input} &
            \textbf{GT} &
            \textbf{Proposed} &
            \textbf{SOGAN} \\
            \includegraphics[width=0.19\linewidth]{../tmbe_deblooming-short-traj-latex-main/images_red/patient_12_vol_10_input_resshift.png} &
            \includegraphics[width=0.19\linewidth]{../tmbe_deblooming-short-traj-latex-main/images_red/patient_12_vol_10_gt_resshift.png} &
            \includegraphics[width=0.19\linewidth]{../tmbe_deblooming-short-traj-latex-main/images_red/patient_12_vol_10_output_resshift.png} &
            \includegraphics[width=0.19\linewidth]{../tmbe_deblooming-short-traj-latex-main/images_red/patient_12_vol_10_output_sogan.png} \\
            \multicolumn{2}{c}{\textbf{RED curve: Proposed}} &
            \multicolumn{2}{c}{\textbf{RED curve: SOGAN}} \\
            \multicolumn{2}{c}{\includegraphics[width=0.34\linewidth]{../tmbe_deblooming-short-traj-latex-main/images_red/patient_12_vol_10_red_curve_resshift.png}} &
            \multicolumn{2}{c}{\includegraphics[width=0.34\linewidth]{../tmbe_deblooming-short-traj-latex-main/images_red/patient_12_vol_10_red_curve_sogan.png}} \\
        \end{tabular}
    \end{minipage}
    \hfill
    \begin{minipage}[t]{0.49\linewidth}
        \centering
        \textbf{Case 2}\\[0.25em]
        \begin{tabular}{cccc}
            \textbf{Input} &
            \textbf{GT} &
            \textbf{Proposed} &
            \textbf{SOGAN} \\
            \includegraphics[width=0.19\linewidth]{../tmbe_deblooming-short-traj-latex-main/images_red/patient_24_vol_2_input_resshift.png} &
            \includegraphics[width=0.19\linewidth]{../tmbe_deblooming-short-traj-latex-main/images_red/patient_24_vol_2_gt_resshift.png} &
            \includegraphics[width=0.19\linewidth]{../tmbe_deblooming-short-traj-latex-main/images_red/patient_24_vol_2_output_resshift.png} &
            \includegraphics[width=0.19\linewidth]{../tmbe_deblooming-short-traj-latex-main/images_red/patient_24_vol_2_output_sogan.png} \\
            \multicolumn{2}{c}{\textbf{RED curve: Proposed}} &
            \multicolumn{2}{c}{\textbf{RED curve: SOGAN}} \\
            \multicolumn{2}{c}{\includegraphics[width=0.34\linewidth]{../tmbe_deblooming-short-traj-latex-main/images_red/patient_24_vol_2_red_curve_resshift.png}} &
            \multicolumn{2}{c}{\includegraphics[width=0.34\linewidth]{../tmbe_deblooming-short-traj-latex-main/images_red/patient_24_vol_2_red_curve_sogan.png}} \\
        \end{tabular}
    \end{minipage}

    \caption{Radial Energy Distribution (RED) analysis on two representative volumes. The images show the central slice of each 3D volume, and the RED curves are computed from the plaque support region.}
    \label{fig:red_analysis}
\end{figure}
\noindent\textbf{Importance of Spatial Prior.}

The first ablation evaluates the role of spatial prior conditioning in the proposed deblooming framework. As shown in Table~\ref{tab:spatial_prior_ablation}, adding the spatial prior substantially improves restoration performance, confirming that localized guidance is important for focusing the model on the blooming-affected region. The qualitative example in Fig.~\ref{fig:mask_ablation} shows the same effect: with spatial prior conditioning, the restored calcium core is more localized and morphologically consistent, while removing the prior leads to errors in both core shape and extent. Notably, the proposed method degrades less sharply without the spatial prior than SOGAN and U-Net++, suggesting that the short-trajectory formulation provides some robustness to missing spatial guidance, although the prior remains important for the best performance.

\begin{table}[t]
\centering
\caption{Ablation study evaluating the effect of calcium spatial prior conditioning. Metrics are aggregated over 232 test volumes from two subjects using fixed HU thresholds: calcium core $>850$ HU and annulus $(550,850]$ HU.}
\label{tab:spatial_prior_ablation}
\resizebox{\linewidth}{!}{
\begin{tabular}{lccc}
\hline
\textbf{Method} &
\textbf{Dice $\uparrow$} &
\textbf{Core RMSE (HU) $\downarrow$} &
\textbf{Annulus RMSE (HU) $\downarrow$} \\
\hline
Blooming-corrupted input & $0.3842 \pm 0.3366$ & $341.4590 \pm 34.8052$ & $176.1679 \pm 14.7582$ \\
\hline
\multicolumn{4}{l}{\textbf{U-Net++}} \\
With spatial prior & $0.8870 \pm 0.0284$ & $187.1475 \pm 27.7500$ & $144.0696 \pm 23.3815$ \\
Without spatial prior & $0.6796 \pm 0.0850$ & $299.7722 \pm 65.0760$ & $289.9809 \pm 68.8215$ \\
\hline
\multicolumn{4}{l}{\textbf{SOGAN}} \\
With spatial prior & $0.8922 \pm 0.0587$ & $131.4178 \pm 19.5148$ & $155.5562 \pm 32.3092$ \\
Without spatial prior & $0.7409 \pm 0.1065$ & $194.8697 \pm 38.0475$ & $232.5836 \pm 32.3222$ \\
\hline
\multicolumn{4}{l}{\textbf{Proposed}} \\
With spatial prior & $0.9250 \pm 0.0781$ & $134.4998 \pm 32.2103$ & $90.2365 \pm 9.7464$ \\
Without spatial prior & $0.8052 \pm 0.2086$ & $236.5011 \pm 49.1845$ & $153.6387 \pm 7.5412$ \\
\hline
\end{tabular}
}
\end{table}

\noindent\textbf{Effect of Soft Annular Boundary Attention}

The second ablation evaluates the effect of Soft Annular Boundary Attention (SABA). We compare the localized ResShift baseline with standard Swin attention in all layers against the proposed version, which uses the same backbone but replaces selected deeper attention blocks with SABA. 
This isolates the effect of the annular boundary penalty beyond the standard Swin-based inversion backbone.

Using an approximate Welch's two-sample $t$-test at $\alpha=0.05$, the proposed version significantly improves Dice ($p=7.12\times10^{-3}$), Core RMSE ($p=6.17\times10^{-24}$), and Annulus RMSE ($p=5.16\times10^{-22}$) over the Swin-only baseline. The qualitative example in Fig.~\ref{fig:saba_ablation} shows the same trend: the Swin-only version recovers a calcium volume close to the ground truth, but the core appears smoother and shows more spreading into the annular region. In contrast, the proposed version better preserves the irregular plaque shape and maintains a cleaner calcium boundary, indicating improved core morphology and calcium-annulus separation.

\begin{table}[t]
\centering
\caption{Ablation study comparing the baseline version with Swin only and the proposed version. Metrics are aggregated over 232 test volumes from two subjects using fixed HU thresholds: calcium core $>850$ HU and annulus $(550,850]$ HU.}
\label{tab:saba_ablation}
\resizebox{\linewidth}{!}{
\begin{tabular}{lccc}
\hline
\textbf{Method} &
\textbf{Dice $\uparrow$} &
\textbf{Core RMSE (HU) $\downarrow$} &
\textbf{Annulus RMSE (HU) $\downarrow$} \\
\hline
Blooming-corrupted input & $0.3842 \pm 0.3366$ & $341.4590 \pm 34.8052$ & $176.1679 \pm 14.7582$ \\
Baseline version with Swin only & $0.9037 \pm 0.1062$ & $165.0378 \pm 29.9030$ & $100.2458 \pm 11.5473$ \\
Proposed version & $0.9250 \pm 0.0781$ & $134.4998 \pm 32.2103$ & $90.2365 \pm 9.7464$ \\
\hline
\end{tabular}
}
\end{table}

\noindent\textbf{Joint Ablation of Spatial Prior and SABA}

\begin{table}[t]
\centering
\caption{Joint ablation of spatial prior and SABA for the proposed model. Metrics are aggregated over 232 test volumes from two subjects using fixed HU thresholds: calcium core $>850$ HU and annulus $(550,850]$ HU.}
\label{tab:combined_ablation}
\resizebox{\linewidth}{!}{
\begin{tabular}{lccc}
\hline
\textbf{Variant} &
\textbf{Dice $\uparrow$} &
\textbf{Core RMSE (HU) $\downarrow$} &
\textbf{Annulus RMSE (HU) $\downarrow$} \\
\hline
Proposed & $0.9250 \pm 0.0781$ & $134.4998 \pm 32.2103$ & $90.2365 \pm 9.7464$ \\
Without SABA & $0.9037 \pm 0.1062$ & $165.0378 \pm 29.9030$ & $100.2458 \pm 11.5473$ \\
Without spatial prior & $0.8052 \pm 0.2086$ & $236.5011 \pm 49.1845$ & $153.6387 \pm 7.5412$ \\
Without spatial prior or SABA & $0.7951 \pm 0.2161$ & $255.2451 \pm 44.8662$ & $160.1108 \pm 14.0153$ \\
\hline
\end{tabular}
}
\end{table}

The joint ablation shows that both components matter: removing either spatial prior conditioning or SABA reduces performance, and removing both causes the largest drop. The full proposed model performs best, which suggests that the two design choices are complementary rather than redundant.

\noindent\textbf{Radial Energy Distribution.}

As a secondary case-based analysis, we examine \textbf{Radial Energy Distribution (RED)}, a metric closely related to the encircled-energy curve used in blur analysis. RED captures how high-intensity calcium signal accumulates with increasing distance from the lesion center in the 3D plaque volume. The full RED setup, including the 550 HU baseline, 850 HU center estimate, and summary radii such as $R_{80}$ and $R_{90}$, is described in Supplementary Sec.~\ref{supp:evaluation_metrics}.

As shown in Fig.~\ref{fig:red_analysis} and Table~\ref{tab:red_r80_r90}, the proposed method stays closer to the GT profile than the blooming-corrupted input, as reflected by the $R_{80}$ and $R_{90}$ values. Relative to the input, it shifts the radial distribution toward the GT reference and reduces the broad blooming tail, consistent with the main qualitative claim that the proposed method better limits halo spread while preserving calcium shape.

\begin{table}[t]
\centering
\caption{$R_{80}$ and $R_{90}$ values for the two RED cases shown in Fig.~\ref{fig:red_analysis}. Smaller deviations from GT indicate a closer match to the reference radial distribution.}
\label{tab:red_r80_r90}
\scriptsize
\setlength{\tabcolsep}{6pt}
\begin{tabular}{lcccc}
\hline
\textbf{Method} & \multicolumn{2}{c}{\textbf{Case 1}} & \multicolumn{2}{c}{\textbf{Case 2}} \\
 & \textbf{$R_{80}$} & \textbf{$R_{90}$} & \textbf{$R_{80}$} & \textbf{$R_{90}$} \\
\hline
Input    & 33.3429 & 35.3664 & 23.0814 & 67.6094 \\
Proposed & 33.4838 & 35.7691 & 25.7171 & 68.2362 \\
SOGAN    & 33.1351 & 35.2712 & 25.3674 & 67.7036 \\
GT       & 34.5010 & 37.0275 & 27.2605 & 69.7661 \\
\hline
\end{tabular}
\end{table}

\section{Discussion}

The results support three main observations. 
First, the baseline comparisons in Table~\ref{tab:deblooming_results} and Figs.~\ref{fig:synthetic_qualitative}--\ref{fig:real_patient_qualitative} indicate that calcium deblooming is not well addressed by generic restoration objectives alone. SRGAN and SAGAN improve apparent sharpness, but remain poorly matched to the calcium-lumen interface, suggesting that visually enhanced boundaries are not sufficient for faithful plaque recovery.  
Second, the results also suggest that deblooming is better modeled as progressive localized restoration than as either single-step correction or long-trajectory diffusion generation. Compared with U-Net++, SOGAN, and the LDM baselines, the proposed method more consistently preserves calcium morphology while suppressing the blooming halo, indicating that a short input-conditioned trajectory is better matched to localized artifact reversal than direct one-step prediction or noise-to-image generation. 
Third, the ablation studies clarify why the full framework is effective. Table~\ref{tab:spatial_prior_ablation} and Fig.~\ref{fig:mask_ablation} indicate that spatial prior conditioning is important for localizing restoration to the blooming-affected region, while Table~\ref{tab:saba_ablation} and Fig.~\ref{fig:saba_ablation} show that SABA improves control at the calcium-lumen boundary beyond the Swin-based backbone alone. Together, these ablations suggest that the proposed method benefits from combining explicit ROI localization with boundary-aware feature control, rather than relying on a single generic restoration mechanism.
Overall, the quantitative metric pattern supports a coherent interpretation of the proposed method. The improved Dice and annulus RMSE indicate stronger plaque-shape recovery and better halo suppression, while the comparable core RMSE suggests that these gains are not achieved by simply eroding the calcium core. Instead, the method appears to better separate true calcium from the surrounding blooming-affected annulus.

This study has several limitations. First, quantitative evaluation is restricted to synthetic paired CCTA data because paired clinical blooming-corrupted and debloomed references are unavailable. Although this allows controlled comparison across methods, it may not fully capture clinical variability in scanner settings, reconstruction kernels, motion, contrast timing, and patient anatomy. Second, the proposed model is trained slice-wise in 2D even though calcium blooming and coronary anatomy are inherently three-dimensional; a fully 3D framework may improve inter-slice consistency and through-plane edge fidelity. Third, clinical CCTA examples are assessed qualitatively only, again due to the lack of paired clinical references. Finally, the current method depends on the availability and quality of the calcium spatial prior. Future work should therefore include larger multi-center clinical studies, reader-based and downstream stenosis-assessment evaluation, and extensions toward jointly learned localization, uncertainty-aware priors, and fully 3D boundary-aware restoration.

\section{Conclusion}

In this paper, we present STAR Diffusion, a short-trajectory diffusion framework for CCTA calcium deblooming. By formulating deblooming as a localized residual-shifting restoration process, the proposed method progressively reduces blooming artifact while preserving calcium structure and supporting lumen recovery. Spatial prior conditioning focuses the restoration on the artifact-affected region, and Soft Annular Boundary Attention (SABA) improves control at the calcium-lumen boundary by reducing interaction with the surrounding blooming annulus. Experiments on synthetic and clinical CCTA data show that this combination produces stronger calcium-region recovery and more consistent blooming suppression than the compared baselines. These results support short-trajectory, spatially guided diffusion as a promising direction for calcium blooming artifact reduction in coronary CT angiography.

\appendix
\section{Reverse Sampling Details}
\label{supp:reverse_sampling}

For the reverse deblooming process, the network predicts the clean target $\hat{x}_0$ at each timestep, and the next sample is drawn from the corresponding Gaussian posterior. Given $\hat{x}_0$, the posterior mean and covariance are
\[
\mu_t
=
\frac{\eta_{t-1}}{\eta_t}x_t
+
\frac{\alpha_t}{\eta_t}\hat{x}_0,
\qquad
\Sigma_t
=
\kappa^2\frac{\eta_{t-1}}{\eta_t}\alpha_t I .
\]
The reverse sample is then obtained as
\[
x_{t-1}
=
\mu_t + \sqrt{\Sigma_t}\,\epsilon,
\qquad \epsilon \sim \mathcal{N}(0,I).
\]
This is the same analytic reverse update used in the main method, written explicitly for clarity.

\section{Radial Energy Distribution}
\label{supp:evaluation_metrics}
Radial Energy Distribution (RED), also known as the encircled-energy curve in optical imaging and blur analysis, measures how high-intensity calcium signal accumulates with distance from the lesion center in the full 3D plaque volume. We define a support region around the calcified plaque using the evaluation mask, estimate the lesion center from the ground-truth high-intensity core ($HU \ge 850$), and convert each voxel to positive radial energy as
\[
E(p)=\max(HU(p)-550,0),
\]
where 550 HU is the baseline so that only elevated attenuation contributing to blooming is accumulated. The cumulative energy is then summed outward as radius increases and normalized by the total energy to produce the RED curve. Summary radii such as $R_{80}$ and $R_{90}$ are used to compare the input, debloomed output, and ground truth; closer agreement with the ground-truth curve indicates better recovery of the radial calcium distribution.

\begin{acks}
This work was supported in part by the National Institutes of Health under Grants EB032716 and HL151561. We thank GE HealthCare Technology and Innovation Center for their support.
\end{acks}

\begin{dci}
The authors declare no conflicts of interest.
\end{dci}

\bibliographystyle{SageV}
\bibliography{zHenriquesLab-Mendeley}

\end{document}
